
\documentclass[11pt]{article}

\usepackage[utf8]{inputenc}
\usepackage[T1]{fontenc}
\usepackage[spanish,es-noshorthands,es-tabla]{babel}
\usepackage[a4paper,margin=2.2cm]{geometry}
\usepackage{amsmath,amssymb,amsfonts,mathtools}
\usepackage{booktabs,longtable,tabularx,array}
\usepackage{ragged2e}
\usepackage{xcolor}
\usepackage{hyperref}
\usepackage{enumitem}
\usepackage{listings}

\hypersetup{
  colorlinks=true,
  linkcolor=blue!55!black,
  urlcolor=blue!55!black,
  citecolor=blue!55!black
}

\setlength{\emergencystretch}{2em}
\renewcommand{\arraystretch}{1.14}
\newcolumntype{L}[1]{>{\RaggedRight\arraybackslash}p{#1}}
\newcolumntype{Y}{>{\RaggedRight\arraybackslash}X}
\DeclareRobustCommand{\codepath}[1]{\begingroup\Urlmuskip=0mu plus 2mu\path{#1}\endgroup}
\pdfstringdefDisableCommands{\def\codepath#1{#1}}

\newcommand{\CaputoD}{{}^{C}D_t^q}
\newcommand{\ii}{\mathrm{i}}
\newcommand{\R}{\mathbb{R}}
\newcommand{\diag}{\operatorname{diag}}
\newcommand{\med}{\operatorname{median}}
\newcommand{\TARGET}{\texttt{TARGET}}

\lstdefinestyle{terminal}{
  basicstyle=\ttfamily\small,
  breaklines=true,
  frame=single,
  columns=fullflexible
}
\lstset{style=terminal}

\title{Complemento matematico y operativo del reporte de la libreria\\
\large \texttt{hidden-attractors-fo}, \texttt{version\_2}}
\author{Revision tecnica del reporte unificado y cruce contra codigo activo}
\date{\today}

\begin{document}
\maketitle
\tableofcontents

\section{Objetivo de este complemento}

El archivo \codepath{version_2/docs/reporte_unificado_chua_fraccionario.tex}
ya contiene el sustento general del proyecto: modelo de Chua fraccionario,
criterio de Matignon, forma de Lur'e, funcion descriptiva, ruta Machado,
verificacion de ocultedad y lectura de figuras. Sin embargo, al cruzarlo contra
los modulos activos de \codepath{version_2/hidden_attractors/}, conviene
agregar varias piezas que si estan codificadas pero no quedan suficientemente
formalizadas en el reporte.

Este documento no reemplaza el reporte principal. Esta escrito como bloque
complementario listo para integrarse despues de las secciones de teoria y antes
de las secciones de figuras/resultados.

\section{Hallazgos de auditoria}

\begin{center}
\begin{tabularx}{\textwidth}{@{}L{0.24\textwidth}L{0.28\textwidth}Y@{}}
\toprule
Elemento revisado & Estado en el reporte actual & Complemento recomendado \\
\midrule
Parametros exactos de la libreria & Aparecen de forma indirecta o en algunos workflows. & Agregar tabla de parametros oficiales usados por \codepath{ChuaParameters}. \\
Equilibrios implementados & Se menciona que se calculan equilibrios, pero falta derivacion algebraica desde el codigo. & Agregar derivacion de \(E_0,E_+,E_-\) con la pendiente \(s_c=-\beta/(\beta+\gamma)\). \\
Jacobiano por regiones & El reporte cita Matignon, pero no fija claramente la matriz por pendiente local. & Agregar \(J(a)\), \(a\in\{m_0,m_1\}\), y el criterio fraccionario. \\
Forma Lur'e efectiva usada por el modelo no suave & El reporte la discute en general. & Agregar \(P,b,r,\psi\) compatibles con la funcion de Chua por tramos. \\
Transferencia fraccionaria & El reporte ya advierte que se usa \(s^q\). & Agregar determinante, cofactores y expresion racional exacta de \(\widehat W_q(\lambda)\). \\
Diferencia entre modelo y contrato numerico & Falta enfatizar que \(q,h,L_m,T,t_{\rm burn}\) no pertenecen al dataclass del modelo. & Agregar separacion: campo vectorial \(F(X)\) versus contrato Caputo/EFORK. \\
Metricas de trayectoria & El reporte menciona FFT, PSD, secciones y nubes, pero no define las formulas usadas. & Agregar definiciones de rango, varianza, entropia espectral, seccion \(x=0\), distancia de nubes. \\
Clases de cuenca & Estan descritas conceptualmente. & Agregar tabla exacta de \codepath{CLASS_LABELS}: equilibrio, target positivo, target negativo, infinito, desconocido, falla numerica. \\
Workflow \codepath{robustness_overlay} & Se menciona robustez. & Agregar casos \(R_0,\ldots,R_5\) y aclarar que no decide ocultedad. \\
Workflow \codepath{sphere_controls} & Se menciona verificacion por esferas. & Agregar construccion \(X_0=E_i+\rho v\), salida \codepath{sphere_decision.csv} y significado de \codepath{hiddenness_status}. \\
Workflow \codepath{refined_basin} & Se menciona refinamiento. & Agregar funcion de puntuacion contra referencias positiva/negativa y umbrales. \\
Adaptadores externos & El reporte lista herramientas. & Agregar que \codepath{nolds} y \codepath{antropy} son adaptadores opcionales, no implementaciones propias. \\
\bottomrule
\end{tabularx}
\end{center}

\section{Sistema implementado: separacion entre modelo y contrato numerico}

La libreria separa dos niveles que en el reporte conviene distinguir
explicitamente:

\begin{enumerate}[label=\textbf{Nivel \arabic*.},leftmargin=2.2cm]
\item \textbf{Modelo algebraico.} Define solamente el campo vectorial
\[
F(X)=
\begin{bmatrix}
\alpha(y-x-f(x))\\
x-y+z\\
-\beta y-\gamma z
\end{bmatrix},
\qquad X=(x,y,z)^T.
\]
\item \textbf{Contrato fraccionario/numerico.} Define como se integra
\[
\CaputoD X=F(X),
\]
con orden \(q\), paso \(h\), memoria finita \(L_m\), tiempo final \(T\),
tiempo de descarte \(t_{\rm burn}\), backend y tolerancias de clasificacion.
\end{enumerate}

Esta separacion es importante porque en el codigo \codepath{ChuaParameters} no
contiene \(q\). El orden fraccionario pertenece a la simulacion Caputo/EFORK y
a los workflows, no al objeto que define el campo vectorial.

\subsection{Parametros oficiales del modelo de Chua no suave}

El dataclass \codepath{ChuaParameters} fija, por defecto,
\[
\alpha=8.4562,\qquad
\beta=12.0732,\qquad
\gamma=0.0052,\qquad
m_0=-0.1768,\qquad
m_1=-1.1468.
\]
La no linealidad implementada es
\begin{equation}
\label{eq:chua_piecewise_impl}
f(x)=m_1x+\frac{m_0-m_1}{2}\bigl(|x+1|-|x-1|\bigr).
\end{equation}
Por regiones:
\[
f(x)=
\begin{cases}
m_1x-(m_0-m_1), & x<-1,\\[2mm]
m_0x, & |x|\leq 1,\\[2mm]
m_1x+(m_0-m_1), & x>1.
\end{cases}
\]

\section{Equilibrios implementados}

Los equilibrios satisfacen \(F(X)=0\). De las dos ultimas ecuaciones:
\[
0=x-y+z,\qquad 0=-\beta y-\gamma z.
\]
De la tercera ecuacion,
\[
z=-\frac{\beta}{\gamma}y,
\]
y al sustituir en la segunda:
\[
x-y-\frac{\beta}{\gamma}y=0
\quad\Longrightarrow\quad
x=\frac{\beta+\gamma}{\gamma}y.
\]
Por tanto,
\begin{equation}
\label{eq:y_z_equilibrium_relation}
y=\frac{\gamma}{\beta+\gamma}x,
\qquad
z=-\frac{\beta}{\beta+\gamma}x.
\end{equation}
La primera ecuacion exige
\[
0=\alpha(y-x-f(x))
\quad\Longrightarrow\quad
f(x)=y-x.
\]
Usando \eqref{eq:y_z_equilibrium_relation},
\[
f(x)=\frac{\gamma}{\beta+\gamma}x-x
=-\frac{\beta}{\beta+\gamma}x.
\]
Definimos la pendiente de cierre
\begin{equation}
\label{eq:closing_slope}
s_c=-\frac{\beta}{\beta+\gamma}.
\end{equation}
Entonces los equilibrios se obtienen resolviendo
\begin{equation}
\label{eq:equilibrium_scalar_condition}
f(x)=s_cx.
\end{equation}

\subsection{Equilibrio central}

En la region \(|x|\leq 1\), \(f(x)=m_0x\). Entonces
\[
m_0x=s_cx
\quad\Longrightarrow\quad
(m_0-s_c)x=0.
\]
Si \(m_0\neq s_c\), se obtiene
\begin{equation}
\label{eq:E0}
E_0=(0,0,0)^T.
\end{equation}

\subsection{Equilibrios externos}

Para \(x>1\),
\[
m_1x+(m_0-m_1)=s_cx,
\]
de donde
\[
(m_1-s_c)x=-(m_0-m_1).
\]
Por tanto,
\begin{equation}
\label{eq:xplus}
x_+=-\frac{m_0-m_1}{m_1-s_c}.
\end{equation}
Para \(x<-1\),
\[
m_1x-(m_0-m_1)=s_cx,
\]
de donde
\[
(m_1-s_c)x=m_0-m_1,
\]
y
\begin{equation}
\label{eq:xminus}
x_-=\frac{m_0-m_1}{m_1-s_c}.
\end{equation}
La funcion auxiliar del codigo construye cada equilibrio como
\[
E(x)=
\begin{bmatrix}
x\\
x+f(x)\\
f(x)
\end{bmatrix}.
\]
Esto coincide con \eqref{eq:y_z_equilibrium_relation}, porque en equilibrio
\(f(x)=s_cx=z\) y \(y=x+f(x)\). Asi,
\begin{equation}
\label{eq:Epm}
E_+=
\begin{bmatrix}
x_+\\
x_+ + f(x_+)\\
f(x_+)
\end{bmatrix},
\qquad
E_-=
\begin{bmatrix}
x_-\\
x_- + f(x_-)\\
f(x_-)
\end{bmatrix}.
\end{equation}

\section{Jacobiano por regiones y criterio fraccionario}

La pendiente local de \(f\) es
\[
a=f'(x)=
\begin{cases}
m_1, & |x|>1,\\
m_0, & |x|<1,
\end{cases}
\]
sin definir en \(x=\pm1\). Lejos de las superficies de quiebre, el Jacobiano es
\begin{equation}
\label{eq:jacobian_region}
J(a)=
\begin{bmatrix}
-\alpha(1+a) & \alpha & 0\\
1 & -1 & 1\\
0 & -\beta & -\gamma
\end{bmatrix},
\qquad a\in\{m_0,m_1\}.
\end{equation}
Para \(E_0\) se usa \(a=m_0\), siempre que \(E_0\) este en la region central.
Para \(E_\pm\) se usa \(a=m_1\), siempre que los puntos externos satisfagan
\(|x_\pm|>1\).

En un sistema fraccionario conmensurado
\[
\CaputoD \xi = J(a)\xi,
\]
la condicion local tipo Matignon es
\begin{equation}
\label{eq:matignon}
|\arg(\lambda_i(J(a)))|>\frac{q\pi}{2}
\qquad
\forall i.
\end{equation}
Si algun autovalor viola \eqref{eq:matignon}, el equilibrio no es localmente
asintoticamente estable bajo ese orden \(q\). Esta prueba es local; no decide
si un atractor observado es oculto.

\section{Forma de Lur'e compatible con el codigo}

A partir de \eqref{eq:chua_piecewise_impl}, se separa
\[
f(x)=m_1x+\psi(x),
\]
donde
\begin{equation}
\label{eq:psi_chua_lure}
\psi(x)=\frac{m_0-m_1}{2}\bigl(|x+1|-|x-1|\bigr).
\end{equation}
Entonces
\[
\CaputoD X = PX+b\,\psi(r^TX),
\]
con
\begin{equation}
\label{eq:Pbr_lure}
P=
\begin{bmatrix}
-\alpha(1+m_1)&\alpha&0\\
1&-1&1\\
0&-\beta&-\gamma
\end{bmatrix},
\qquad
b=
\begin{bmatrix}
-\alpha\\0\\0
\end{bmatrix},
\qquad
r=
\begin{bmatrix}
1\\0\\0
\end{bmatrix}.
\end{equation}
La salida escalar de la no linealidad es
\[
\sigma=r^T X=x.
\]

\section{Transferencia fraccionaria y calculo explicito}

La transferencia del bloque lineal fraccionario se escribe
\begin{equation}
\label{eq:Wq}
W_q(s)=r^T(s^qI-P)^{-1}b.
\end{equation}
Para evitar confundir el eje de frecuencia entero con el fraccionario, se
introduce
\[
\lambda=s^q,
\qquad
\widehat W_q(\lambda)=r^T(\lambda I-P)^{-1}b.
\]
En frecuencia:
\begin{equation}
\label{eq:lambda_fractional_frequency}
\lambda=(\ii\omega)^q
=\omega^q\left(
\cos\frac{q\pi}{2}
+\ii\sin\frac{q\pi}{2}
\right).
\end{equation}

Con \eqref{eq:Pbr_lure},
\[
\lambda I-P=
\begin{bmatrix}
\lambda+\alpha(1+m_1)&-\alpha&0\\
-1&\lambda+1&-1\\
0&\beta&\lambda+\gamma
\end{bmatrix}.
\]
Definimos
\begin{equation}
\label{eq:Nlambda}
N(\lambda)=(\lambda+1)(\lambda+\gamma)+\beta
=\lambda^2+(1+\gamma)\lambda+(\beta+\gamma).
\end{equation}
El determinante se obtiene expandiendo por la primera fila:
\begin{align}
\Delta(\lambda)
&=\det(\lambda I-P)\\
&=
\bigl(\lambda+\alpha(1+m_1)\bigr)
\det
\begin{bmatrix}
\lambda+1&-1\\
\beta&\lambda+\gamma
\end{bmatrix}
-(-\alpha)
\det
\begin{bmatrix}
-1&-1\\
0&\lambda+\gamma
\end{bmatrix}\\
&=
\bigl(\lambda+\alpha(1+m_1)\bigr)N(\lambda)
-\alpha(\lambda+\gamma).
\end{align}
Como \(b=(-\alpha,0,0)^T\), se necesita la entrada \((1,1)\) de
\((\lambda I-P)^{-1}\). Por cofactores,
\[
[(\lambda I-P)^{-1}]_{11}
=
\frac{C_{11}}{\Delta(\lambda)}
=
\frac{N(\lambda)}{\Delta(\lambda)}.
\]
Por tanto,
\begin{equation}
\label{eq:W_hat_explicit}
\widehat W_q(\lambda)
=
-\alpha\,
\frac{N(\lambda)}{\Delta(\lambda)}.
\end{equation}
Con la convencion \(\widehat W_q=r^T(\lambda I-P)^{-1}b\) y
\(b=(-\alpha,0,0)^T\), la condicion armonica de primer orden se interpreta como
\begin{equation}
\label{eq:harmonic_condition}
1-N_{\rm DF}(A)\,
\widehat W_q\!\left((\ii\omega)^q\right)=0,
\end{equation}
donde \(N_{\rm DF}(A)\) es una funcion descriptiva aproximada. En el uso del
proyecto, \eqref{eq:harmonic_condition} genera semillas; no prueba ciclos
periodicos exactos del sistema causal de Caputo.

\section{Puente Weyl--Caputo para justificar el uso armonico}

El balance armonico y la funcion descriptiva trabajan naturalmente con señales
periodicas estacionarias. En calculo fraccionario esto corresponde mejor al
marco de Weyl/Liouville--Weyl que al problema causal de Caputo iniciado en
\(t_0\). La lectura usada por la libreria debe formularse asi:

\begin{enumerate}[label=\textbf{(\roman*)},leftmargin=1.6cm]
\item Caputo modela la dinamica causal con memoria desde \(t_0\).
\item Weyl/Liouville--Weyl permite representar respuestas periodicas ideales
sobre historia infinita.
\item La discrepancia Caputo--Weyl aparece como un termino de memoria
dependiente de la historia inicial.
\item Por tanto, el calculo armonico produce una semilla asintotica o
aproximada para explorar el sistema Caputo, no una solucion periodica exacta
garantizada.
\end{enumerate}

Este punto debe agregarse cerca de la seccion de funcion descriptiva para
evitar una conclusion mas fuerte que la evidencia.

\section{Funcion descriptiva de Machado como generador auxiliar de semillas}

La extension tipo Machado debe documentarse como una deformacion compleja de
una funcion descriptiva clasica:
\begin{equation}
\label{eq:machado_fdf_general}
N_{\star,\mu}(A)=\exp\{\mu\,\operatorname{Log}N_\star(A)\},
\qquad \mu>0,
\end{equation}
donde \(\star\) identifica la no linealidad base y \(\operatorname{Log}\) debe
fijar una rama compleja. En particular:
\[
\mu=1
\quad\Longrightarrow\quad
N_{\star,\mu}(A)=N_\star(A).
\]
Valores \(\mu\neq1\) no cambian el orden de Caputo \(q\). Solo deforman la
ganancia compleja usada para proponer semillas. En consecuencia, los candidatos
Machado deben pasar por las mismas pruebas EFORK, robustez y cuencas que los
candidatos Lur'e clasicos.

\section{Metricas de trayectoria implementadas}

Las trayectorias se leen bajo la convencion
\[
\texttt{t,x,y,z}.
\]
Despues de un tiempo de analisis \(t_{\rm start}\), se usa la cola
\[
\mathcal T_{\rm tail}
=
\{X(t_k):t_k\geq t_{\rm start}\}.
\]

\subsection{Rangos y varianzas}

Para \(u\in\{x,y,z\}\),
\begin{equation}
\label{eq:range_u}
R_u=\max_{t_k\geq t_{\rm start}} u(t_k)
-\min_{t_k\geq t_{\rm start}} u(t_k).
\end{equation}
Tambien se calcula
\[
\operatorname{var}_u
=
\operatorname{Var}\{u(t_k):t_k\geq t_{\rm start}\}.
\]
Estas cantidades detectan colapso a equilibrio o dinamica no trivial, pero no
son pruebas de caos.

\subsection{Frecuencia dominante y entropia espectral}

La funcion \codepath{component_fft} centra la senal, aplica ventana de Hann y
calcula potencia espectral:
\[
\tilde u_n=w_n(u_n-\bar u),
\qquad
P_k=|\mathcal F(\tilde u)_k|^2.
\]
La frecuencia dominante ignora la componente \(k=0\):
\[
k_\ast=\arg\max_{k\geq1}P_k,
\qquad
f_\ast=f_{k_\ast}.
\]
La entropia espectral normalizada se calcula como
\begin{equation}
\label{eq:spectral_entropy}
H_{\rm PSD}
=
-\frac{1}{\log M}
\sum_{k=1}^{M}p_k\log(p_k+10^{-300}),
\qquad
p_k=\frac{P_k}{\sum_{\ell=1}^{M}P_\ell}.
\end{equation}
Esta metrica indica dispersion espectral; no reemplaza el calculo de exponentes
de Lyapunov.

\subsection{Seccion de Poincare operacional}

La libreria usa cruces ascendentes de \(x=0\). Para dos puntos consecutivos con
\[
x_{k-1}<0\leq x_k,
\]
se exige ademas
\[
F_1(X_k)>0.
\]
El punto \((y,z)\) de la seccion se interpola linealmente:
\[
\eta=\frac{0-x_{k-1}}{x_k-x_{k-1}},
\]
\[
y_\Sigma=y_{k-1}+\eta(y_k-y_{k-1}),
\qquad
z_\Sigma=z_{k-1}+\eta(z_k-z_{k-1}).
\]
La nube de seccion es
\[
\Sigma=\{(y_\Sigma,z_\Sigma)\}.
\]

\subsection{Distancia simetrica entre nubes}

Para comparar dos nubes \(A\) y \(B\), el codigo usa una distancia mediana de
vecino mas cercano:
\begin{equation}
\label{eq:cloud_median_distance}
d_{\rm cloud}(A,B)
=
\med\left(
\left\{
\min_{b\in B}\|a-b\|:a\in A
\right\}
\cup
\left\{
\min_{a\in A}\|b-a\|:b\in B
\right\}
\right).
\end{equation}
Se usa para comparar geometria post-transitoria, no para probar igualdad
punto a punto entre trayectorias caoticas.

\section{Clases de cuenca y lectura operacional}

La clasificacion de cuencas usa etiquetas enteras. La tabla exacta que debe
aparecer en el reporte es:

\begin{center}
\begin{tabularx}{\textwidth}{@{}cL{0.24\textwidth}Y@{}}
\toprule
ID & Etiqueta & Lectura operacional \\
\midrule
0 & \texttt{equilibrium} & La trayectoria cae o queda cerca de un equilibrio.\\
1 & \texttt{target\_positive} & Clase target positiva, bounded no trivial.\\
2 & \texttt{target\_negative} & Clase target negativa, bounded no trivial.\\
3 & \texttt{infinity} & Divergencia o salida de cota.\\
4 & \texttt{unknown} & No clasificada bajo el contrato usado.\\
5 & \texttt{numerical\_failure} & Falla numerica o error del backend.\\
6 & \texttt{bounded\_other} & Solo en refinamiento: bounded no colapsada, pero no cercana al target.\\
\bottomrule
\end{tabularx}
\end{center}

Los IDs target son
\[
\mathcal C_{\rm target}=\{1,2\}.
\]
Un \codepath{target_hit} significa
\[
\operatorname{class\_id}(X_0)\in\{1,2\}.
\]
Esto no prueba que el atractor sea oculto; al contrario, si ocurre desde una
vecindad de algun equilibrio, es evidencia contra la ocultedad bajo el contrato
probado.

\section{Workflow \texttt{robustness\_overlay}}

El workflow \codepath{robustness_overlay} compara trayectorias iniciadas desde
el mismo candidato bajo variaciones de \(h\), \(L_m\) y \(T\). La familia
estandar es:

\begin{center}
\begin{tabular}{lcccc}
\toprule
Caso & \(q\) & \(h\) & \(L_m\) & \(T\) \\
\midrule
\(R_0\) base & \(q\) & 0.01 & 10 & 1500\\
\(R_1\) paso fino & \(q\) & 0.005 & 10 & 1500\\
\(R_2\) paso grueso & \(q\) & 0.02 & 10 & 1500\\
\(R_3\) memoria menor & \(q\) & 0.01 & 5 & 1500\\
\(R_4\) memoria mayor & \(q\) & 0.01 & 20 & 1500\\
\(R_5\) tiempo largo & \(q\) & 0.01 & 10 & 3000\\
\bottomrule
\end{tabular}
\end{center}

Su salida principal es \codepath{robustness_overlay_metrics.csv}.
La lectura correcta es:
\[
\text{geometria persistente bajo perturbaciones}
\not\Longrightarrow
\text{atractor oculto}.
\]
Solo indica que el atractor observado no desaparece facilmente al modificar el
contrato numerico probado.

\section{Workflow \texttt{sphere\_controls}}

El workflow \codepath{sphere_controls} prueba condiciones iniciales sobre
esferas alrededor de cada equilibrio:
\begin{equation}
\label{eq:sphere_initial_condition}
X_0=E_i+\rho v,
\qquad
\|v\|=1,
\qquad
E_i\in\{E_0,E_+,E_-\}.
\end{equation}
Los radios por defecto son
\[
\rho\in\{10^{-5},3\cdot10^{-5},10^{-4},3\cdot10^{-4},10^{-3}\}.
\]
El resultado se resume en \codepath{sphere_decision.csv}. La regla de decision
codificada es:
\begin{itemize}
\item si \codepath{total_target_hits} es positivo, entonces \codepath{hiddenness_status=not_supported_by_sphere_equilibrium_test};
\item si \codepath{total_target_hits} es cero, entonces \codepath{hiddenness_status=compatible_with_hiddenness_under_tested_spheres}.
\end{itemize}

La segunda salida no significa ``oculto demostrado''. Solo significa que la
muestra de esferas y radios no encontro contacto con la cuenca target.

\section{Workflow \texttt{refined\_basin}}

El workflow \codepath{refined_basin} revisa celdas previamente etiquetadas como
\codepath{unknown}. Para cada celda reintegra el punto inicial y compara la
geometria con atractores de referencia positivo y negativo.

Para una trayectoria candidata y una referencia \(R\), calcula:
\[
d_{\rm cloud}^{\rm norm},\qquad
d_{\Sigma}^{\rm norm},\qquad
d_{\rm range},\qquad
d_{\rm FFT}.
\]
El puntaje se puede escribir como
\begin{equation}
\label{eq:refined_score}
S_R=
\frac{
w_c d_{\rm cloud}^{\rm norm}
+w_r d_{\rm range}
+w_f \min(d_{\rm FFT},2)
+w_s d_{\Sigma}^{\rm norm}
}{
w_c+w_r+w_f+w_s
},
\end{equation}
usando los terminos disponibles. Por defecto:
\[
w_c=1,\qquad w_r=0.35,\qquad w_f=0.20,\qquad w_s=0.35.
\]
La referencia ganadora es
\[
R_\ast=\arg\min_{R\in\{+,-\}}S_R.
\]
Una celda se reclasifica como target si satisface simultaneamente los umbrales
del contrato:
\[
S_{R_\ast}\leq0.95,\qquad
d_{\rm cloud}^{\rm norm}\leq0.85,\qquad
d_{\rm range}\leq1.50,\qquad
d_{\rm FFT}\leq1.00,\qquad
d_\Sigma^{\rm norm}\leq1.20.
\]
Si es bounded y no colapsada pero no cumple esos umbrales, recibe clase
\[
\texttt{bounded\_other}.
\]

\section{Candidatos: que representa realmente un \texttt{CandidateRecord}}

Un registro candidato almacena informacion de semilla, ruta y punto final de
continuacion. Matematicamente, debe leerse como
\[
\mathcal R=
(\text{id},\text{ruta},q,X_{\rm seed},X_{\rm robust},A,\sigma_0,\omega,\rho_H,\ldots).
\]
Los campos principales son:

\begin{center}
\begin{tabular}{ll}
\toprule
Campo & Interpretacion \\
\midrule
\codepath{candidate_id} & Identificador reproducible del candidato.\\
\codepath{route} & Ruta de origen: Lur'e, Machado/FDF, etc.\\
\codepath{q} & Orden fraccionario usado en la corrida asociada.\\
\codepath{seed} & Semilla inicial propuesta por funcion descriptiva/continuacion.\\
\codepath{robust_start} & Punto desde el cual se hacen pruebas de robustez.\\
\codepath{A}, \codepath{omega} & Amplitud y frecuencia armonica aproximadas, si existen.\\
\codepath{rho_H} & Diagnostico heuristico de alejamiento/contacto, si existe.\\
\codepath{source} & Archivo de salida desde donde se cargo el registro.\\
\bottomrule
\end{tabular}
\end{center}

Un \codepath{CandidateRecord} no es evidencia de ocultedad. Es un paquete de
metadatos para lanzar verificaciones posteriores.

\section{Adaptadores externos y alcance}

Los modulos de integracion externa no implementan desde cero algoritmos de
complejidad. Delegan a paquetes opcionales como \codepath{nolds} y
\codepath{antropy}.
Esto implica que cualquier resultado de entropia, dimension de correlacion,
DFA, Hurst o Lyapunov tipo Rosenstein debe citar tambien el paquete externo y
su metodo, no solo la libreria local.

Para bifurcaciones, la funcion local
\codepath{bifurcation_points_from_trajectories} solo postprocesa trayectorias ya calculadas. No hace continuacion numerica ni
seguimiento de ramas. Si se usa continuacion real, debe documentarse la
herramienta externa correspondiente.

\section{Regla de lectura final}

La ruta completa que debe defenderse en el reporte es:
\[
\boxed{
\begin{gathered}
\text{semilla}
\to
\text{integracion Caputo/EFORK}
\to
\text{atractor bounded no trivial}
\to
\text{robustez}\\
\to
\text{cuencas}
\to
\text{pruebas desde equilibrios}
\to
\text{lectura de ocultedad}
\end{gathered}
}
\]
La funcion descriptiva y la ruta Machado pertenecen al primer bloque:
\[
\text{generacion de semillas}.
\]
La decision de ocultedad pertenece al ultimo bloque:
\[
\text{cuencas alrededor de todos los equilibrios}.
\]
Por tanto, en el reporte no debe escribirse:
\begin{quote}
``La funcion descriptiva encontro un atractor oculto''.
\end{quote}
Debe escribirse:
\begin{quote}
``La funcion descriptiva genero una semilla que, bajo el contrato numerico probado, conduce a un atractor candidato; la ocultedad se evalua mediante cuencas''.
\end{quote}

\section{Texto breve para insertar como advertencia metodologica}

\begin{quote}
Todas las etiquetas producidas por los workflows son etiquetas numericas bajo
un contrato finito de simulacion. En particular,
\codepath{compatible\_with\_hiddenness\_under\_tested\_spheres} no significa
``atractor oculto demostrado'', sino ``no se encontro interseccion con las
vecindades de equilibrio muestreadas bajo los radios, tiempos, memoria,
tolerancias y backend usados''. La conclusion debe mantenerse condicionada al
contrato numerico y a la resolucion de la cuenca explorada.
\end{quote}

\section{Bibliografia adicional recomendada para completar el reporte}

\begin{thebibliography}{99}

\bibitem{Machado2014}
J. Tenreiro Machado.
\emph{Fractional order describing functions}.
Signal Processing, 2014.
DOI: 10.1016/j.sigpro.2014.05.012.

\bibitem{LeonovKuznetsov2013}
G. A. Leonov and N. V. Kuznetsov.
\emph{Hidden Attractors in Dynamical Systems: From Hidden Oscillations in
Hilbert--Kolmogorov, Aizerman, and Kalman Problems to Hidden Chaotic Attractor
in Chua Circuits}.
International Journal of Bifurcation and Chaos, 2013.

\bibitem{Kuznetsov2016}
N. V. Kuznetsov.
\emph{Hidden Attractors in Dynamical Systems: Systems with No Equilibria,
Multistability and Coexisting Attractors}.
Physics Reports, 2016.

\bibitem{Matignon1996}
D. Matignon.
\emph{Stability Results for Fractional Differential Equations with Applications
to Control Processing}.
Computational Engineering in Systems Applications, 1996.

\bibitem{DiethelmFordFreed2002}
K. Diethelm, N. J. Ford, and A. D. Freed.
\emph{A Predictor-Corrector Approach for the Numerical Solution of Fractional
Differential Equations}.
Nonlinear Dynamics, 2002.

\bibitem{TavazoeiHaeri2009}
M. S. Tavazoei and M. Haeri.
\emph{A Proof for Non Existence of Periodic Solutions in Time Invariant
Fractional Order Systems}.
Automatica, 2009.

\end{thebibliography}

\end{document}
